% =====================================================================
%  CHƯƠNG 6. PHÂN LOẠI PHƯƠNG TIỆN                | Phụ trách: Nhật
%  Viết high-level, không nhắc mã nguồn/hàm/đường dẫn.
%  Model triển khai: loại xe = ResNet18; kiểu dáng = MobileNetV3-Small.
% =====================================================================
\chapter{Phân Loại Phương Tiện}
\label{ch:phanloai}

Module phân loại phương tiện gồm ba bước nối tiếp: (1) \textbf{loại xe} (mô hình tự huấn
luyện, chạy trên cả khung hình, cho nhãn ô tô, xe máy hoặc xe tải); (2) \textbf{định vị xe}
(mô hình YOLO pretrained COCO, không huấn luyện thêm, chỉ tìm và cắt vùng xe cho bước~3);
(3) \textbf{kiểu dáng xe con} (chỉ chạy khi bước~1 ra ô tô và bước~2 cắt được vùng xe; mô
hình tự huấn luyện, ba lớp Sedan, Gầm cao, Xe tải). Bước~1 thay cho việc dùng
thẳng nhãn lớp COCO của YOLO, vốn sai nhiều trên ảnh cận cảnh camera cổng
(Mục~\ref{sec:phanloai-ketqua}).

\textbf{Về phạm vi so với đề cương.} Đề cương yêu cầu bốn lớp (ô tô, xe máy, xe tải, xe
buýt); mô hình hiện phân loại được ba lớp. Lớp xe buýt chưa có mô hình riêng do thiếu ảnh xe
khách đúng góc camera cổng, và nhãn xe buýt của YOLO pretrained bị bỏ vì đo được là sai.

\section{Dữ Liệu}
\label{sec:phanloai-dulieu}

\subsection{Dữ liệu cho bước loại xe}

Nguồn là bộ A1 (cùng dữ liệu module phát hiện biển số), chia theo năm nhóm ảnh với nội dung
khác nhau, đã xem tay xác nhận (Bảng~\ref{tab:vt-sources}).

\begin{table}[htbp]
\centering
\caption{Các nhóm ảnh trong bộ A1 dùng cho bước loại xe.}
\label{tab:vt-sources}

\begin{tabular}{lrp{7.5cm}}
\toprule
\textbf{Nhóm} & \textbf{Số ảnh} & \textbf{Nội dung} \\
\midrule
Ô tô cận cảnh & 989 & Ô tô kiểu camera barrier, thuần một loại \\
Xe máy cận cảnh & 1.747 & Xe máy kiểu CCTV cổng, thuần một loại \\
Ảnh phố hỗn hợp & 925 & Nhiều loại xe, hậu cảnh lộn xộn \\
Cắt cực sát & 917 & Chỉ thấy biển và ca-lăng, không dùng cho phân loại loại xe \\
\bottomrule
\end{tabular}
\end{table}

\textbf{Lần huấn luyện đầu mắc lỗi học tắt (shortcut learning).} Khi mỗi lớp chỉ lấy từ đúng
một nguồn ảnh (ô tô từ ảnh cận cảnh, xe máy từ ảnh cận cảnh, xe tải từ ảnh phố), mô hình đạt
100\% trên tập test riêng nhưng chỉ đúng 1 trên 13 ảnh phố thật ngoài tập train, gần như mọi
ảnh bị đoán thành xe tải, vì đó là lớp duy nhất có ảnh phong cách đường phố. Mô hình đã học
nguồn ảnh và góc camera thay vì hình dáng xe.

Cách khắc phục là trộn phong cách ảnh cho từng lớp, buộc mô hình học hình dáng xe: bổ sung
ảnh phố cho lớp ô tô và xe máy bên cạnh ảnh cận cảnh gốc, và xem tay toàn bộ ứng viên cho lớp
xe tải vì độ tin cậy của detector không đủ để lọc tự động. Tập chia 70/15/15 với seed cố định
(Bảng~\ref{tab:vt-split}); lớp xe tải mỏng hơn nhiều nên mất cân bằng được xử lý bằng trọng
số lớp trong hàm mất mát.

\begin{table}[htbp]
\centering
\caption{Số ảnh theo lớp và split của bước loại xe.}
\label{tab:vt-split}

\begin{tabular}{lrrrr}
\toprule
\textbf{Lớp} & \textbf{Train} & \textbf{Valid} & \textbf{Test} & \textbf{Tổng} \\
\midrule
Ô tô & 309 & 66 & 67 & 442 \\
Xe máy & 335 & 72 & 71 & 478 \\
Xe tải & 51 & 11 & 11 & 73 \\
\midrule
\textbf{Tổng} & 695 & 149 & 149 & 993 \\
\bottomrule
\end{tabular}
\end{table}

\subsection{Dữ liệu cho bước kiểu dáng xe con}

Bộ B5 (Vehicle Body Style Dataset, Roboflow, CC BY 4.0) gồm 10.000 ảnh (train 7.014, valid
1.999, test 987) với 12 kiểu dáng gốc, là nguồn duy nhất cho bước kiểu dáng. Mười hai kiểu
dáng được gộp thành ba nhóm thực dụng cho nhu cầu vận hành bãi xe (Bảng~\ref{tab:b5-mapping}).
Ảnh được cắt theo bounding box gốc với biên 10\%, giữ nguyên split
(Bảng~\ref{tab:b5-split}, Hình~\ref{fig:phanloai-phanboo}).

\begin{table}[htbp]
\centering
\caption{Gộp 12 kiểu dáng gốc B5 thành 3 lớp phân loại.}
\label{tab:b5-mapping}
{\footnotesize Nguồn: bộ B5~\autocite{researchprojects2022vehiclebodystyle}, gộp lớp theo đề tài.}
\begin{tabular}{ll}
\toprule
\textbf{Lớp mới} & \textbf{Kiểu dáng gốc B5 gộp vào} \\
\midrule
Sedan & Sedan, Fastback, Hatchback, Wagon, Convertible, Hardtop Convertible, Sports \\
Gầm cao & SUV, Crossover, MPV, Minibus \\
Xe tải & Pickup Truck \\
\bottomrule
\end{tabular}
\end{table}

\begin{table}[htbp]
\centering
\caption{Số ảnh theo lớp và split sau khi gộp từ B5.}
\label{tab:b5-split}

\begin{tabular}{lrrr}
\toprule
\textbf{Lớp} & \textbf{Train} & \textbf{Valid} & \textbf{Test} \\
\midrule
Sedan & 4.102 & 1.166 & 569 \\
Gầm cao & 2.325 & 670 & 334 \\
Xe tải & 587 & 163 & 84 \\
\midrule
\textbf{Tổng} & 7.014 & 1.999 & 987 \\
\bottomrule
\end{tabular}
\end{table}

\begin{figure}[htbp]
\centering
\includegraphics[width=0.90\textwidth]{vehicle_style_class_distribution.png}
\caption{Phân bố lớp theo split trong bộ dữ liệu B5 sau khi gộp.}
\label{fig:phanloai-phanboo}

\end{figure}

\section{Kiến Trúc và Huấn Luyện}
\label{sec:phanloai-kientruc}

Cả hai bước phân loại so sánh hai kiến trúc fine-tune từ trọng số ImageNet, thay lớp cuối cho
ba lớp: ResNet18~\autocite{he2016resnet} (khoảng 11,2 triệu tham số) và
MobileNetV3-Small~\autocite{howard2019mobilenetv3} (khoảng 1,5 triệu tham số, nhẹ hơn khoảng
7,3 lần). Điểm khác biệt là đầu vào: bước loại xe dùng \textbf{cả khung hình} (resize
224$\times$224) vì dữ liệu huấn luyện là ảnh camera cổng nguyên khung, còn bước kiểu dáng
dùng \textbf{vùng xe đã cắt} vì huấn luyện trên ảnh B5 cắt theo bounding box. Bước định vị
dùng YOLOv8n pretrained COCO~\autocite{jocher2023yolov8}, lấy bounding box xe lớn nhất kèm
biên 10\%.

Cả hai bước dùng chung cấu hình huấn luyện (Bảng~\ref{tab:app-cls}), mỗi mô hình dưới 7 phút
trên GPU; đường hội tụ nhanh và ổn định (Hình~\ref{fig:phanloai-losscurves}).

\begin{figure}[htbp]
\centering
\includegraphics[width=0.90\textwidth]{vehicle_style_loss_curves.png}
\caption{Đường mất mát huấn luyện và accuracy kiểm định theo epoch của hai kiến trúc (bước kiểu dáng).}
\label{fig:phanloai-losscurves}

\end{figure}

\section{Kết Quả}
\label{sec:phanloai-ketqua}

\subsection{Bước loại xe}

Đánh giá trên tập test 149 ảnh (67 ô tô, 71 xe máy, 11 xe tải) sau khi khắc phục lỗi học tắt
(Bảng~\ref{tab:vt-results}, chi tiết từng lớp ở Phụ lục~\ref{app:ketqualop}). Cả hai mô hình
gần như tuyệt đối với ô tô và xe máy; xe tải yếu nhất (9/11) vì ít ảnh nhất, và lỗi phân tán
đều sang hai lớp kia, không còn dấu hiệu học tắt.

\begin{table}[htbp]
\centering
\caption{So sánh ResNet18 và MobileNetV3-Small ở bước loại xe (tập test 149 ảnh).}
\label{tab:vt-results}
{\footnotesize CPU inference đo trên máy huấn luyện, không phải Raspberry Pi~5.}
\begin{tabular}{lrrrrr}
\toprule
\textbf{Model} & \textbf{Accuracy} & \textbf{F1-macro} & \textbf{Params} & \textbf{ONNX} & \textbf{CPU inference} \\
\midrule
ResNet18 & 0,9866 & 0,9619 & 11,2 tr & 44,7 MB & 8,19 ms \\
MobileNetV3-Small & 0,9799 & 0,9571 & 1,5 tr & 6,1 MB & 4,22 ms \\
\bottomrule
\end{tabular}
\end{table}

Để xác nhận lỗi học tắt đã khắc phục tận gốc, đề tài dựng một bộ kiểm tra cố định gồm 12 ảnh
thật ngoài mọi tập train, valid và test: ResNet18 đúng 12/12 (100\%), MobileNetV3-Small 11/12
(92\%), so với chỉ 1/13 ở lần huấn luyện mắc lỗi học tắt.

Accuracy trên đây là của riêng mô hình. Khi ghép vào pipeline, nếu đặt bước định vị YOLO làm
điều kiện chạy thì detector pretrained bỏ sót 57\% ảnh cận cảnh camera cổng. Vì vậy đề tài cho mô hình loại xe chạy trên cả khung hình (đúng như dữ liệu
huấn luyện), chỉ dùng bằng chứng có xe (thấy xe hoặc đọc được biển) làm điều kiện an toàn;
nhờ đó tỉ lệ ảnh có kết quả đạt 100\% và độ chính xác end-to-end đạt 98,7\%, khớp đúng
accuracy của mô hình. Xe khách, chưa có lớp riêng, hiện do nhân viên sửa tay.

\subsection{Bước kiểu dáng}

Đánh giá trên tập test B5 987 ảnh (Bảng~\ref{tab:phanloai-sosanh}, Hình~\ref{fig:phanloai-sosanh}).
So với bài toán 12 lớp thử trước đó, accuracy cải thiện rõ (ResNet18 từ 0,768 lên 0,900;
MobileNetV3-Small từ 0,701 lên 0,901); với ba lớp, MobileNetV3-Small gần như ngang ResNet18
dù nhẹ hơn 7,3 lần. Ma trận nhầm lẫn (Hình~\ref{fig:phanloai-confusion}) cho thấy Sedan và Xe
tải phân loại tốt (93,1\% đến 96,4\% tùy mô hình); nhầm lẫn tập trung ở Gầm cao (79,6\% với
ResNet18, 83,2\% với MobileNetV3-Small), chủ yếu bị đoán nhầm thành Sedan, hợp lý vì crossover
cỡ nhỏ đôi khi có dáng gần sedan.

\begin{table}[htbp]
\centering
\caption{So sánh ResNet18 và MobileNetV3-Small ở bước kiểu dáng (tập test B5, 987 ảnh).}
\label{tab:phanloai-sosanh}
{\footnotesize CPU inference đo trên máy huấn luyện, không phải Raspberry Pi~5.}
\begin{tabular}{lrrrrr}
\toprule
\textbf{Model} & \textbf{Accuracy} & \textbf{F1-macro} & \textbf{Params} & \textbf{ONNX} & \textbf{CPU inference} \\
\midrule
ResNet18 & 0,8997 & 0,9012 & 11,2 tr & 44,7 MB & 8,57 ms \\
MobileNetV3-Small & 0,9007 & 0,9069 & 1,5 tr & 6,1 MB & 4,82 ms \\
\bottomrule
\end{tabular}
\end{table}

\begin{figure}[htbp]
\centering
\includegraphics[width=0.90\textwidth]{vehicle_style_comparison.png}
\caption{So sánh tổng quan hai kiến trúc theo các chỉ số đánh giá (bước kiểu dáng).}
\label{fig:phanloai-sosanh}

\end{figure}

\begin{figure}[htbp]
\centering
\includegraphics[width=0.90\textwidth]{vehicle_style_confusion_matrices.png}
\caption{Ma trận nhầm lẫn trên tập test của ResNet18 và MobileNetV3-Small (bước kiểu dáng).}
\label{fig:phanloai-confusion}

\end{figure}

\section{Thảo Luận}
\label{sec:phanloai-thaoluan}

\textbf{Accuracy trên tập test nội bộ không đảm bảo hiệu năng thật.} Kết quả 1/13 trước và
12/12 sau khi trộn nguồn trên bộ ngoài phân phối cho thấy tập kiểm tra phải tách nguồn với
tập huấn luyện; độ đa dạng nguồn ảnh quan trọng hơn số lượng ảnh.

\textbf{Accuracy của mô hình không phải accuracy người dùng nhận được.} Khi ghép nhiều mô
hình, không nên để một thành phần chưa đánh giá (bước định vị YOLO, bỏ sót 57\% ảnh) chặn
đầu vào của thành phần đã có số liệu.

\textbf{Hạn chế còn lại.} YOLO pretrained COCO vẫn bỏ sót xe trên nhiều ảnh camera cổng; khi
đó hệ thống vẫn ra loại xe và biển số nhưng thiếu khung xe và kiểu dáng. Mô hình kiểu dáng
huấn luyện hoàn toàn trên ảnh đại lý của B5 nên con số 90\% chưa được kiểm chứng trên ảnh
camera cổng thật.

\textbf{Chọn mô hình triển khai.} Bước loại xe chọn ResNet18 vì chính xác hơn rõ (0,9866 so
với 0,9799 trên tập test, 100\% so với 92\% trên bộ ngoài phân phối). Bước kiểu dáng chọn
MobileNetV3-Small: hai kiến trúc gần ngang nhau về accuracy (0,8997 so với 0,9007) nhưng
MobileNetV3-Small nhẹ hơn 7,3 lần, giảm footprint và chạy mát hơn mà không giảm độ chính xác
đáng kể (đo trên Raspberry Pi~5 ở Chương~\ref{ch:trienkhai}).
